\section{Methodology}

\subsection{State-Space Model Specification}

Our structural state-space model for inflation decomposes the inflation rate into unobserved components while incorporating multivariate economic features.

\subsubsection{General Framework}

The state-space representation consists of two equations:

\begin{align}
y_t &= \mathbf{Z}_t \boldsymbol{\alpha}_t + \boldsymbol{\beta}' \mathbf{x}_t + \varepsilon_t, \quad \varepsilon_t \sim N(0, H_t) \label{eq:obs}\\
\boldsymbol{\alpha}_t &= \mathbf{T}_t \boldsymbol{\alpha}_{t-1} + \mathbf{R}_t \boldsymbol{\eta}_t, \quad \boldsymbol{\eta}_t \sim N(0, \mathbf{Q}_t) \label{eq:state}
\end{align}

where:
\begin{itemize}
    \item $y_t$ is the observed inflation rate at time $t$
    \item $\boldsymbol{\alpha}_t$ is the $(m \times 1)$ state vector containing unobserved components
    \item $\mathbf{x}_t$ is the $(k \times 1)$ vector of exogenous features
    \item $\boldsymbol{\beta}$ is the $(k \times 1)$ vector of feature coefficients
    \item $\mathbf{Z}_t$, $\mathbf{T}_t$, $\mathbf{R}_t$ are system matrices
    \item $\varepsilon_t$ and $\boldsymbol{\eta}_t$ are observation and state innovations
\end{itemize}

\subsubsection{Component Decomposition}

The state vector $\boldsymbol{\alpha}_t$ contains:

\begin{enumerate}
    \item \textbf{Trend component} ($\mu_t$, $\nu_t$): Long-run inflation level and slope
    \begin{align}
        \mu_t &= \mu_{t-1} + \nu_{t-1} + \eta_{\mu,t}, \quad \eta_{\mu,t} \sim N(0, \sigma^2_\mu)\\
        \nu_t &= \nu_{t-1} + \eta_{\nu,t}, \quad \eta_{\nu,t} \sim N(0, \sigma^2_\nu)
    \end{align}
    
    \item \textbf{Cyclical component} ($\gamma_t$): Business cycle fluctuations
    \begin{align}
        \gamma_t = \rho \gamma_{t-1} + \eta_{\gamma,t}, \quad \eta_{\gamma,t} \sim N(0, \sigma^2_\gamma)
    \end{align}
    where $|\rho| < 1$ ensures stationarity.
    
    \item \textbf{Seasonal component} (optional for monthly/quarterly data)
\end{enumerate}

\subsubsection{Multivariate Features}

The feature vector $\mathbf{x}_t$ includes:
\begin{itemize}
    \item Macroeconomic indicators: GDP growth, unemployment rate, industrial production
    \item Monetary variables: Money supply growth, interest rates
    \item External factors: Exchange rates, commodity prices
    \item Expectations: Survey-based inflation expectations
\end{itemize}

\subsection{Approximate Bayesian Computation}

\subsubsection{ABC Framework}

Given observed data $y_{1:T}^{obs}$, we seek the posterior distribution:
\begin{align}
    p(\boldsymbol{\theta} | y_{1:T}^{obs}) \propto p(y_{1:T}^{obs} | \boldsymbol{\theta}) p(\boldsymbol{\theta})
\end{align}

where $\boldsymbol{\theta}$ contains all model parameters (variance parameters, AR coefficients, feature weights).

When the likelihood $p(y_{1:T}^{obs} | \boldsymbol{\theta})$ is intractable, ABC approximates the posterior by:

\begin{enumerate}
    \item Sample $\boldsymbol{\theta}^*$ from prior $p(\boldsymbol{\theta})$
    \item Simulate data $y_{1:T}^{sim}$ from the model with $\boldsymbol{\theta}^*$
    \item Compute summary statistics $\mathbf{s}(y_{1:T})$
    \item Accept $\boldsymbol{\theta}^*$ if $d(\mathbf{s}(y^{obs}), \mathbf{s}(y^{sim})) < \varepsilon$
\end{enumerate}

\subsubsection{Summary Statistics}

We use summary statistics that capture key features of inflation:
\begin{itemize}
    \item Mean and variance
    \item Autocorrelations at various lags
    \item Quantiles
    \item Volatility measures
\end{itemize}

\subsubsection{ABC-SMC Algorithm}

To improve efficiency, we implement Sequential Monte Carlo ABC:

\begin{algorithm}[H]
    \caption{ABC-SMC for State-Space Models}
    \begin{algorithmic}
        \FOR{$t = 1$ to $T$}
            \STATE Set tolerance $\varepsilon_t$ (decreasing sequence)
            \FOR{$i = 1$ to $N$}
                \IF{$t = 1$}
                    \STATE Sample $\boldsymbol{\theta}^{(i)}$ from prior
                \ELSE
                    \STATE Sample from previous particles with weights
                    \STATE Perturb using kernel
                \ENDIF
                \STATE Simulate data with $\boldsymbol{\theta}^{(i)}$
                \STATE Compute distance $d_i$
                \IF{$d_i < \varepsilon_t$}
                    \STATE Accept and compute weight
                \ENDIF
            \ENDFOR
        \ENDFOR
    \end{algorithmic}
\end{algorithm}

\subsection{Forecasting}

Point forecasts are generated via:
\begin{align}
    \hat{y}_{T+h|T} = E[y_{T+h} | y_{1:T}]
\end{align}

Prediction intervals use the posterior predictive distribution:
\begin{align}
    p(y_{T+h} | y_{1:T}) = \int p(y_{T+h} | \boldsymbol{\theta}) p(\boldsymbol{\theta} | y_{1:T}) d\boldsymbol{\theta}
\end{align}
